\documentclass[11pt]{article}
\usepackage[margin=27mm]{geometry}
\usepackage{amsmath,amssymb,amsthm,booktabs,longtable,hyperref}
\newtheorem{theorem}{Theorem}
\newtheorem{proposition}[theorem]{Proposition}
\newtheorem{corollary}[theorem]{Corollary}
\newcommand{\Sz}{\operatorname{Sz}}
\newcommand{\Aut}{\operatorname{Aut}}
\newcommand{\F}{\mathbb F}
\title{A corrected classification of rank-two Suzuki geometries}
\author{Human Compression Project}
\date{}
\begin{document}
\maketitle
\begin{abstract}
We assemble a complete treatment of the rank-two classification in
Leemans's 1998 paper. The solvable parabolic families follow from a common
normalizer incidence argument. The missing subfield families are the
minimal central covers of a subfield point stabilizer, parametrized by
projective lines over the smaller field. Their Frobenius orbits give a
corrected classification and exact counting formulas. Integer certificates
settle all sixty rows of the numerical table; literal permutation and
coset certificates settle the twenty-nine diagrams for $\Sz(8)$.
The printed statements and their corrections remain distinct. Classical
Suzuki structure is used as an attributed premise.
\end{abstract}

\section{Scope and conventions}
The original is D.~Leemans, \emph{The rank 2 geometries of the simple
Suzuki groups $\Sz(q)$}, Beitr\"age zur Algebra und Geometrie \textbf{39}
(1998), 97--120. We retain its equivalence convention: an automorphism
of the specified group sends one subgroup pair to the other, with type
interchange permitted. Inner equivalence uses inner automorphisms.
Thus the classification concerns a geometry together with its acting
group. It does not assert that all underlying unmarked incidence graphs
are nonisomorphic.

Throughout, $G=\Sz(q)$ with $q=2^n$, $n\ge3$ odd. A configuration is an
unordered pair of proper subgroups $M,L$ such that $J=M\cap L$ is maximal
in each, and at least one of $M,L$ is maximal in $G$. The corresponding
coset geometry is firm, residually connected, flag transitive and
residually weakly primitive. The following elementary observation supplies
the rank-two reduction used in source Lemma~2.1.

\begin{proposition}
Suppose $M$ is maximal in $G$, $L\nleq M$, and $J=M\cap L$ is maximal
in both parabolics. Then the coset geometry on $G/M$ and $G/L$, with
incidence defined by nonempty intersection, has the four properties above.
\end{proposition}
\begin{proof}
Every incident pair has a common representative $g$, and hence is the
translate of the base pair $(M,L)$. This proves flag transitivity.
The two local degrees are $[M:J]$ and $[L:J]$, both at least two, giving
firmness. Its incidence graph is connected exactly when
$\langle M,L\rangle=G$: a path is an alternating sequence of cosets
and therefore gives a word in the two subgroups; conversely each such
word gives a path. Here generation follows from maximality of $M$.
Rank-one residues have no further connectivity condition. The action
on $G/M$ is primitive by maximality, and both residue actions are
primitive by maximality of $J$ in the corresponding parabolic. These
are precisely the residual weak-primitivity conditions in rank two.
\end{proof}

The classical premises are the faithful Suzuki root coordinates, the
ovoid action and involution centralizers, the torus normalizers, the
subgroup and maximal-subgroup classifications, and
$\Aut(G)=G\rtimes\operatorname{Gal}(\F_q/\F_2)$. These are attributed to
Suzuki, Tits and L\"uneburg, as in the original Sections~2.2 and~4 and
Lemma~6.7. In particular, Tits's ovoid and normal-form theorem in source
Lemma~2.3 is retained as a premise, not presented as a new proof.
The elementary field assertion in Lemma~2.2 follows because the
finite-field automorphisms are the powers of Frobenius: the equation
$\sigma^2:x\mapsto x^2$ is $2i\equiv1\pmod n$, with exactly one solution
for odd $n$ and none for even $n$.

\section{The common arguments and the complete families}
We use the detailed proofs in the accompanying notes [2--5]. The first
reproves the arithmetic preliminaries and the maximal subgroups of the
solvable normalizers, and excludes a point stabilizer as a parabolic.
One affine-group argument treats source Lemmas~5.1--5.3. The exclusion of
a point stabilizer also gives the original assertion in every rank greater than one of
Lemma~6.1; that assertion is not inferred merely from the rank-two case.

The normalizer argument counts embedded overgroups by counting incidences
between two subgroup conjugacy classes. Fixing the intersection leaves
a normalizer action; inversion and Frobenius then account for type
exchange and outer automorphisms. In root coordinates the relevant
finite sets are quotients of a Suzuki two-group or its centre. This
recovers all solvable maximal-parabolic families and their orbit counts
in source Lemmas~6.2--6.6 and~6.8, with the explicit Table~6 correction
recorded below.

For the remaining maximal parabolic, write $n=dr$, where $r$ is prime,
$d\ge3$, $K=\F_{2^d}$ and $F=\F_{2^n}$. Let $S=\Sz(2^d)$ and let
$B_s$ be its point stabilizer, where $s=2^d$. In the faithful root
coordinates write $Z_T=\{u(0,t):t\in T\}$. The crucial replacement for
the printed subfield row is
\[
 Q_T=Z_TB_s,\qquad K\le T\le F,\qquad\dim_KT=2.       \tag{1}
\]
These groups have order $s^3(s-1)$. The source instead gives the larger
group of order $qs(s-1)$; that group does not cover $B_s$ when $r\ge3$.

\begin{theorem}[Complete corrected classification]\label{classification}
All configurations are the solvable maximal-parabolic families of [4],
together with the following families having maximal parabolic $S$:
\[
\begin{array}{c|c|l}
J=S\cap L&L&\text{condition}\\\hline
 B_s&Q_T&\text{as in (1)}\\
 D_{2(s-1)}&D_{2m}&m\mid q-1,\quad m/(s-1)\text{ prime}\\
 C_{a_\delta(s)}\rtimes C_4&C_m\rtimes C_4&
 m\mid a_\epsilon(q),\quad m/a_\delta(s)\text{ prime}\\
 \Sz(2^e)&\Sz(2^{er})&n=erp,\quad r\ne p\text{ prime},\ e\ge3.
\end{array}                                                   \tag{2}
\]
Here $a_\pm(x)=x\pm\sqrt{2x}+1$, and $\epsilon$ is the unique ambient
sign containing $a_\delta(s)$. Central-cover classes are the Frobenius
orbits on $\mathbb P_K(F/K)$. Each permitted numerical choice in the
second or third row gives one class. The last row gives one class for
each unordered pair of distinct prime indices. Families appearing
under both maximal parabolics are identified once.
\end{theorem}
\begin{proof}
The classical maximal-subgroup list gives a point stabilizer, the three
torus normalizers, and the maximal simple subfield groups. The first is
excluded by [2], and [4] treats the three normalizers completely.
For $S$, its maximal-subgroup list leaves the intersections in (2).

The key exhaustion argument for the first row is worth isolating.
Any solvable overgroup of $B_s$ lies in its uniquely determined ambient
point stabilizer $B_q=P\rtimes F^*$. Put $V=L\cap P$.
If $B_s$ is maximal in $L$, the normalizer condition in the two-group
$V$ forces its subfield root subgroup $U_K$ to be normal in $L$.
The root-coordinate calculation in [5] gives
$N_P(U_K)=Z_FU_K$ and $N_{B_q}(U_K)=Z_FU_KK^*$.
Consequently $L\le Z_FB_s$. The subgroup interval above $B_s$ is
exactly the lattice of $K$-subspaces containing $K$, by [3]. Its
minimal elements are (1). Thus these are all covers, not only examples
omitted from the printed table. Their equivalence is the stated
semilinear line equivalence.

The characteristic odd torus in either of the next two intersections
forces its unique ambient torus normalizer. Maximality then makes the
torus extension a prime extension. A simple overgroup having that
intersection maximal would be the original $S$, by the uniqueness of
a subfield overgroup of prescribed order. For a simple intersection
$\Sz(2^e)$, maximality in each parabolic gives two prime relative degrees;
uniqueness excludes equal indices. The case $2^e=2$ is solvable and
belongs to the preceding nonsplit row, so it is not counted here.
These are the exhaustion arguments proved in full in [5]. Finally,
type interchange accounts for the stated identifications.
\end{proof}

\section{The exact counting formulas}
Let $\omega(m)$ be the number of distinct prime divisors of $m$. For the
first row of (2), define $C(n,d)$ as follows. For $0\le j<n$, put
$h=\gcd(d,j)$ and $g=\gcd(n,j)$, using $\gcd(a,0)=a$. Then
\[
 C(n,d)=\frac1n\sum_{j=0}^{n-1}b_j,\qquad
 b_j=\begin{cases}
 (2^{h(r-1)}-1)/(2^h-1),&g=rh,\\
 \gcd(r,2^h-1)-1,&g=h.
 \end{cases}                                               \tag{3}
\]
There are no other cases. The proof in [5] counts fixed projective lines
by finite-field Hilbert 90 and averaging; it also covers $r\mid d$.
Thus no coprimality of the absolute and relative degrees is assumed.

For composite $n$, the total number $T(n)$ of configurations is
\[
 T(n)=1+E(n)+\sum_{\substack{r\mid n\ \mathrm{prime}\\d=n/r\ge3}}
                \bigl(C(n,d)+A(n,d)\bigr),                  \tag{4}
\]
where
\begin{align*}
 E(n)&=\#\{\{r,p\}:r\ne p\text{ prime},\ rp\mid n,\ n/(rp)\ge3\},\\
 A(n,d)&=\omega\!\left(\frac{q-1}{2^d-1}\right)
       +\sum_{\delta=\pm1}\omega\!\left(
              \frac{t_\delta(q,2^d)}{a_\delta(2^d)}\right).
\end{align*}
Here $t_\delta(q,s)$ denotes the unique $a_\epsilon(q)$ divisible by
$a_\delta(s)$. The constant term is the universal pair
$(D_{2(q-1)},\F_q\rtimes\F_q^*)$.

For prime $n$, put
\[
 v=\frac{q-2}{2n},\quad I=1+\frac{(q-2)(q+4)}{4n},\quad
 H=1+\frac{(q-2)(q+2)}{4n}.
\]
Let $\eta_0,\eta_+,\eta_-$ indicate primality of $q-1,a_+(q),a_-(q)$.
Then
\begin{align}
T(n)={}&1+\eta_0[I+2v+1+H(\omega(a_+(q))+\omega(a_-(q)))]\notag\\
&+\sum_{\epsilon=\pm1}\eta_\epsilon
 [2v+(v+1)\omega(a_{-\epsilon}(q))]
 -\eta_+\eta_-(v+1).                                    \tag{5}
\end{align}
Both formulas follow by adding the complete family counts. In (4), the
simple subfield pair is counted by an unordered pair of indices; in
(5), the last term removes the twice-counted mixed nonsplit family.
These are the only overlaps, as proved in [5]. In particular the source
Corollary~7.1 is recovered: if $n$ is prime and all three odd factors are
composite, then $T(n)=1$.

\section{Numerical and diagram applications}
The integer-only criterion in [6] certifies all three odd factors for
each intended exponent $e=1,\ldots,60$, where $n=2e+1$. The result is
$180$ complete factorizations using $694$ recursively certified primes.
The exact products, primality witnesses, prime counts and all sixty
substitutions in (4)--(5) are retained as readable data and certificates.
Of the printed total entries, eighteen exact values and eleven lower
bounds are recovered. Thirty exact values require correction, and the
lower bound at $e=39$ is false:
\[
 T(79)=19128573095168183144089<10^{23}.
\]
The minus factor has three distinct prime divisors at $e=43$ and four
at the intended $e=60$, rather than the printed two and three.
The last row is literally labelled $30$; its position and parameters
specify $60$. We record that contextual correction explicitly.

For $q=8$, [7] acquires the already verified permutation model and
subgroup classes. Incidence counting proves exhaustion of the embedded
pairs; normalizer actions and type exchange give their equivalence
classes. Equivariant label maps certify the actual coset spaces, and
literal graph searches certify diameters, girths and local double
transitivity. There are exactly $29$ outer classes and $75$ inner
classes, recovering Corollary~8.1. All printed diagram entries are joined
to these data, subject to the following explicit corrections:
\begin{itemize}
\item In diagrams 25 and 26, residue orders $(12,4)$ become $(4,12)$
 in the printed parabolic orientation.
\item Their inner multiplicities are interchanged: the edge triple
 $(6,4,6)$ has multiplicity one and $(5,3,6)$ has multiplicity three.
\item Diagrams 27 and 29 are not residually quasiprimitive. The order-eight
 parabolic admits three different saturated chains to $G$ [8].
\end{itemize}
The other positive primitivity marks are proved in [7], not inferred
from graph appearance. Repeated indistinguishable diagram entries are
compared as multisets; no arbitrary unique identification is imposed.

\section{Disposition of the source and attribution}
The arithmetic and subgroup preliminaries in Lemmas~3.1--3.7 and
5.1--5.3 are recovered by [2]. Lemma~3.2 needs its proper-subfield domain
made explicit: $q=s^r$, $s>2$, $r>1$. Lemma~6.1 is recovered in all ranks greater than one.
The solvable families in Lemmas~6.2--6.6 and~6.8, and Tables~2--7, follow
from [4]. In Table~6 the condition that $p$ divide $\beta_q$ is replaced
by $p\mid\alpha_q$, while the hypothesis that $\beta_q$ be prime remains.
Lemma~6.7 and Table~1 are the attributed classical subgroup inputs.

Lemma~6.9 and Table~8 require the complete replacement (1). The corrected
Theorem~7.1 and Tables~9--11 are Theorem~\ref{classification} and
(3)--(5), with the local double-transitivity classification of [5].
The universal pair is locally doubly transitive, as is every central
cover. The only additional locally doubly transitive families occur at
$q=8$, with parabolics $(C_5\rtimes C_4,C_5\rtimes C_4)$ or
$(C_5\rtimes C_4,C_4\times C_2)$. Therefore the Section~9 assertion that
only the universal family survives for $e>2$ is false when $n$ is
composite. For prime $n>5$ it is recovered. The quoted numerical examples
at $e=3,6,9$ are respectively $2411$, $6453727$ and $18084183211$.
The informal assertion that subfield cases have only a small number of
geometries supplies no formal bound; the corrected exact totals replace
the printed numerical evidence for it.

The source's announced future studies of higher ranks and of
$\operatorname{PSL}(2,q)$ are not asserted results of this manuscript.
They are not converted into extra obligations or solved-conjecture claims.
Every stated classification family, numerical row and rank-two diagram
of the fixed original has a disposition above. The complete treatment
retains the original refutation of the subfield row and supplies proofs
of the corrected classification and all surviving claims.

The contribution of this treatment is the common normalizer and scalar
subspace arguments, their complete exhaustion and orbit analysis, and the
explicit arithmetic and graph certificates. The source results remain
attributed to their original authors; classical Suzuki structure and
standard finite-field and group-action results remain borrowed premises.
No claim of priority for the standard tools, no Lean verification, and no
measured compression ratio is made here. The present assembly is submitted
for distinct whole-manuscript review.

\begin{thebibliography}{9}
\bibitem{source} D. Leemans, \emph{The rank 2 geometries of the simple Suzuki
groups $\Sz(q)$}, Beitr\"age zur Algebra und Geometrie \textbf{39} (1998),
97--120.
\bibitem{pre} Human Compression Project, \emph{Maximal subgroups and elementary reductions
for Suzuki geometries}, accompanying proof note.
\bibitem{scalar} Human Compression Project, \emph{Scalar subspaces in rank-two Suzuki
geometries}, accompanying proof note.
\bibitem{normal} Human Compression Project, \emph{Normalizer fibres in
rank-two Suzuki geometries}, accompanying proof note.
\bibitem{subfield} Human Compression Project, \emph{Subfield parabolics and
the corrected Suzuki classification}, accompanying proof note.
\bibitem{arith} Human Compression Project, \emph{Exact arithmetic for the
Suzuki classification table}, accompanying proof note.
\bibitem{diagram} Human Compression Project, \emph{Finite certificates for
the rank-two geometries of $\Sz(8)$}, accompanying proof note.
\bibitem{labels} Human Compression Project, \emph{Two corrections to the diagrams for $\Sz(8)$}, accompanying proof note.
\end{thebibliography}
\end{document}
